\chapter{El algoritmo \emph{fake news}}
\label{cap:fakenews}

Este es el capítulo central. Contiene el resultado técnico de \ABRS\ y la receta
que se deriva de él. Vale la pena leerlo despacio: la demostración es corta y
entenderla cambia por completo la sensación de estar usando una caja negra.

\section{El punto de partida: el método directo}

Supongamos que queremos la columna $s$ del jacobiano. El procedimiento obvio es:

\begin{enumerate}[leftmargin=1.4em,itemsep=1pt]
\item Perturbar el insumo en la fecha $s$: $X^s = X_\sst + e_s\, dx$.
\item Iterar \eqref{eq:canon-v} hacia atrás desde $T-1$ hasta $0$, guardando
$v^s_t$, $\Lambda^s_t$ y $y^s_t$ para cada $t$.
\item Iterar \eqref{eq:canon-D} hacia adelante desde $D_0 = D_\sst$.
\item Agregar: $Y^s_t = (y^s_t)' D^s_t$, y formar $\J_{t,s} = (Y^s_t - Y_\sst)/dx$.
\end{enumerate}

Hay que repetirlo $T$ veces, una por columna. Con $T=300$ eso son $300$ pasadas
hacia atrás y $300$ hacia adelante, cada una de $300$ pasos: unos $180\,000$
pasos en total. En el modelo HANK de dos activos de \abrs\ eso toma 956 segundos.
El algoritmo fake news lo hace en 3.5 segundos.

\section{Primer resultado: las políticas solo dependen de la distancia}

\begin{lema}[Invariancia de las políticas]
\label{lema:politicas}
Sean $y^s_t$ y $\Lambda^s_t$ la función de resultados y la matriz de transición
en la fecha $t$ ante un choque en la fecha $s$. Entonces
\begin{equation}
y^s_t = \begin{cases} y^{T-1}_{T-1-(s-t)} & s\ge t\\ y_\sst & s<t\end{cases}
\qquad
\Lambda^s_t = \begin{cases} \Lambda^{T-1}_{T-1-(s-t)} & s\ge t\\ \Lambda_\sst & s<t\end{cases}
\label{eq:lema1}
\end{equation}
\end{lema}

La demostración es inmediata a partir de la estructura recursiva de
\eqref{eq:canon-v}: $v_t$ depende de $v_{t+1}$ y $X_t$, y nada más. Si el choque
está $s-t$ períodos en el futuro, la solución hacia atrás es idéntica sea cual
sea el valor de $t$ en el calendario. Y si el choque ya pasó ($s<t$), el agente
se comporta como en el estado estacionario, porque \eqref{eq:canon-v} no tiene
memoria del pasado.

\begin{idea}
El lema \ref{lema:politicas} vale \textbf{sin aproximación}: es cierto para
choques de cualquier tamaño, no solo infinitesimales. Y ya basta por sí solo
para eliminar $T-1$ de las $T$ pasadas hacia atrás: una sola iteración hacia
atrás, a partir de un choque en la fecha $T-1$, entrega las políticas para todas
las fechas de choque y todas las fechas de respuesta.
\end{idea}

Queda el problema de las $T$ pasadas hacia adelante. Para eso hace falta el
segundo resultado, que sí requiere infinitesimalidad.

\section{Segundo resultado: la noticia falsa}

Definamos, para $s\ge1$ y $t\ge1$,
\begin{equation}
\F_{t,s}\cdot dx \;\equiv\; dY^s_t - dY^{s-1}_{t-1}.
\label{eq:def-F}
\end{equation}
Esto compara la respuesta en $t$ a un choque en $s$ con la respuesta un período
antes a un choque un período antes. Por el lema \ref{lema:politicas}, las
políticas en ambos casos son idénticas. Uno podría conjeturar que la diferencia
es cero. No lo es, porque la \emph{distribución} también cambió.

\begin{lema}[Diferencia de respuestas]
\label{lema:fake}
Si $D_0 = D_\sst$, entonces para $dx$ infinitesimal y $s,t\ge1$,
\begin{equation}
\F_{t,s}\cdot dx = y_\sst'\,(\Lambda_\sst')^{\,t-1}\, dD^s_1 .
\label{eq:lema2}
\end{equation}
\end{lema}

\begin{proof}
Linealizando \eqref{eq:canon-Y} alrededor del estado estacionario,
$dY^s_t = y_\sst'\,dD^s_t + (dy^s_t)'\,D_\sst$. Restando la misma expresión en
$(t-1,s-1)$ y usando $dy^s_t = dy^{s-1}_{t-1}$ del lema \ref{lema:politicas},
el segundo término se cancela y queda
\begin{equation}
\F_{t,s}\cdot dx = y_\sst'\big(dD^s_t - dD^{s-1}_{t-1}\big).
\end{equation}
Ahora linealizamos \eqref{eq:canon-D}:
$dD^s_t = \Lambda_\sst'\,dD^s_{t-1} + (d\Lambda^s_{t-1})'D_\sst$. Restando su
análoga en $(t-1,s-1)$ y usando otra vez el lema \ref{lema:politicas}
($d\Lambda^s_{t-1} = d\Lambda^{s-1}_{t-2}$), el segundo término se cancela y
queda la recursión
\begin{equation}
dD^s_t - dD^{s-1}_{t-1} = \Lambda_\sst'\big(dD^s_{t-1} - dD^{s-1}_{t-2}\big).
\end{equation}
Iterando hasta la base, y usando $dD^{s-1}_0 = 0$ (porque $D_0 = D_\sst$ está
dado), se obtiene $dD^s_t - dD^{s-1}_{t-1} = (\Lambda_\sst')^{t-1} dD^s_1$,
que es \eqref{eq:lema2}.
\end{proof}

\paragraph{De dónde sale el nombre} La expresión \eqref{eq:lema2} tiene una
lectura económica preciosa. $\F_{t,s}$ es la respuesta agregada a un choque
particular: \emph{se anuncia en la fecha $0$ que el insumo cambiará en la fecha
$s$, y en la fecha $1$ el anuncio se retracta}. Los agentes reaccionan al
anuncio, lo que mueve la distribución de $D_\sst$ a $D^s_1$. Después, al saber
que era falso, vuelven a comportarse como en el estado estacionario. Entonces el
efecto sobre los agregados en toda fecha $t\ge1$ es simplemente el de haber
partido de una distribución distinta, propagado con la dinámica estacionaria.
De ahí: \emph{noticia falsa}.

\section{Los vectores de expectativa}

La expresión \eqref{eq:lema2} se simplifica con una definición.

\begin{definicion}[Vector de expectativa]
\label{def:expectativa}
$\E_t \equiv (\Lambda_\sst)^t\, y_\sst$.
\end{definicion}

Para cada punto de la grilla, $\E_t$ dice cuál es el valor esperado del resultado
$Y$ dentro de $t$ períodos para un agente que hoy está en ese punto, bajo
comportamiento de estado estacionario. Con esto, $\F_{t,s}\cdot dx = \E_{t-1}'\,dD^s_1$,
y los $\E_t$ se obtienen de una recursión trivial, $\E_t = \Lambda_\sst \E_{t-1}$
con $\E_0 = y_\sst$, que no depende del choque en absoluto.

\begin{figure}[htbp]
\centering
\includegraphics[width=\textwidth]{fig04_vectores_expectativa.pdf}
\caption{Izquierda: vectores de expectativa del ahorro agregado para el estado
de ingreso intermedio. Un hogar con muchos activos hoy seguirá teniendo muchos
activos dentro de cuatro trimestres, pero la relación se aplana a medida que
$t$ crece. Derecha: la dispersión de $\E_t$ entre puntos de la grilla, que mide
cuánta información sobre el agregado futuro contiene la posición actual en la
distribución. Decae, pero despacio.}
\label{fig:expectativa}
\end{figure}

El panel derecho de la figura \ref{fig:expectativa} tiene una lectura económica
directa: cuantifica cuánto tiempo ``recuerda'' la economía una perturbación de
la distribución. En esta calibración, después de 160 trimestres todavía queda
cerca de un tercio de la dispersión inicial. Esa persistencia es exactamente la
razón por la que $T$ tiene que ser grande.

\section{El resultado principal}

\begin{proposicion}[\ABRS, proposición 1]
\label{prop:principal}
Supóngase $D_0 = D_\sst$. Defínase la \emph{matriz fake news} $\F$ por
\begin{equation}
\F_{t,s}\cdot dx =
\begin{cases}
dY^s_0 & t=0\\[2pt]
\E_{t-1}'\, dD^s_1 & t\ge1
\end{cases}
\label{eq:propF}
\end{equation}
donde $dY^s_0 = (dy^s_0)'D_\sst$ y $dD^s_1 = (d\Lambda^s_0)'D_\sst$. Entonces el
jacobiano satisface la recursión
\begin{equation}
\J_{t,s} = \J_{t-1,s-1} + \F_{t,s},
\qquad\text{es decir}\qquad
\J_{t,s} = \sum_{k=0}^{\min\{s,t\}} \F_{t-k,\,s-k}.
\label{eq:recursion}
\end{equation}
\end{proposicion}

La recursión \eqref{eq:recursion} dice que cada entrada del jacobiano es la suma
de la entrada correspondiente de $\F$ y de todas las entradas de $\F$ sobre la
diagonal que va hacia arriba y a la izquierda. Por ejemplo,
$\J_{2,2} = \F_{2,2} + \F_{1,1} + \F_{0,0}$.

\begin{figure}[htbp]
\centering
\includegraphics[width=\textwidth]{fig03_matriz_fakenews.pdf}
\caption{La matriz fake news del ahorro agregado respecto de la tasa. La primera
columna coincide con la del jacobiano. Las columnas posteriores son el efecto de
una noticia anunciada y retractada: se desacumula desde el primer período,
porque lo único que queda del anuncio es la distribución perturbada que dejó.}
\label{fig:fakenews}
\end{figure}

Comparando la figura \ref{fig:fakenews} con la \ref{fig:columnas} se ve la
economía de la descomposición. Las columnas de $\F$ son monótonas y decaen; las
de $\J$ tienen forma de carpa. La carpa aparece porque \eqref{eq:recursion}
acumula columnas desplazadas de $\F$: la acumulación previa al choque es la suma
de muchos términos pequeños de ``noticias falsas'' sucesivas.

\section{La receta}

\begin{receta}[Algoritmo fake news: un insumo y un producto]
\begin{enumerate}[leftmargin=1.4em,itemsep=3pt]
\item \textbf{Una iteración hacia atrás.} Perturbar el insumo y resolver
\eqref{eq:canon-v} hacia atrás una sola vez. De ahí salen, para cada
$s=0,\dots,T-1$:
\begin{itemize}[leftmargin=1.2em,topsep=2pt]
\item los $T$ escalares $\mathcal{Y}_s\,dx \equiv dY^s_0 = (dy^s_0)'D_\sst$:
el efecto sobre el producto \emph{hoy} de una noticia sobre la fecha $s$;
\item los $T$ vectores $\mathcal{D}_s\,dx \equiv dD^s_1 = (d\Lambda^s_0)'D_\sst$:
el efecto sobre la distribución \emph{de mañana} de esa misma noticia.
\end{itemize}
\item \textbf{Los vectores de expectativa.} $\E_0 = y_\sst$ y
$\E_t = \Lambda_\sst \E_{t-1}$ para $t=1,\dots,T-2$. No dependen del choque.
\item \textbf{Armar $\F$.} Primera fila: los $\mathcal{Y}_s$. Filas restantes:
los productos internos $\E_{t-1}'\mathcal{D}_s$. En código es un único producto
de matrices de tamaño $(T-1)\times n_g$ por $n_g \times T$.
\item \textbf{Acumular.} $\J_{t,s} = \J_{t-1,s-1}+\F_{t,s}$, recorriendo filas.
\end{enumerate}
\end{receta}

\paragraph{Varios insumos y productos} Los $\mathcal{D}_s$ dependen solo del
insumo: se calculan una vez por insumo. Los $\E_t$ dependen solo del producto:
una vez por producto. Solo los $\mathcal{Y}_s$ dependen de ambos. Por eso un
bloque con cinco insumos y cuatro productos --- veinte jacobianos --- cuesta
mucho menos que veinte veces un jacobiano.

\paragraph{Por qué es tan rápido} El método directo hace $T$ pasadas hacia atrás
y $T$ hacia adelante, cada una de $T$ pasos. El algoritmo fake news hace una
pasada hacia atrás de $T$ pasos y $T$ aplicaciones de $\Lambda_\sst$. La mejora
es de un factor $\approx T$, lo que con $T=300$ significa dos órdenes de
magnitud.


\section{El algoritmo con números: un recorrido completo}

Nada aclara tanto como ver las matrices. Esta sección ejecuta el algoritmo sobre
el bloque de hogares de Krusell--Smith con $T=6$, un horizonte absurdamente corto
pero que cabe en la página. Todas las cifras salen del código del apéndice
\ref{ap:codigo}.

\paragraph{La matriz fake news} El paso 3 produce
\begin{equation}
\F^{A,r} =
\begin{pmatrix}
35.834 & 0.636 & 0.611 & 0.589 & 0.568 & 0.548 \\
35.550 & 0.623 & 0.600 & 0.579 & 0.558 & 0.539 \\
35.268 & 0.611 & 0.589 & 0.569 & 0.549 & 0.530 \\
34.987 & 0.601 & 0.579 & 0.559 & 0.540 & 0.522 \\
34.709 & 0.590 & 0.570 & 0.550 & 0.531 & 0.514 \\
34.432 & 0.580 & 0.560 & 0.541 & 0.523 & 0.506
\end{pmatrix}
\end{equation}

Vale la pena leer su estructura antes de seguir.

\begin{itemize}[leftmargin=1.4em]
\item \textbf{La primera columna es enorme} comparada con las demás (35.8 frente
a 0.6). Es el efecto riqueza de un choque \emph{contemporáneo}: cuando la tasa
sube hoy, los hogares reciben un retorno extra sobre todo el acervo de activos
--- que es grande --- y lo ahorran casi todo.
\item \textbf{Las demás columnas son pequeñas y muy parecidas entre sí}. Son el
efecto de una noticia anunciada y retractada: los hogares ahorran un poco en el
período 0 anticipando el choque, y eso es todo lo que queda.
\item \textbf{Cada columna decae suavemente hacia abajo}. Ese decaimiento es la
disipación de la perturbación distributiva, es decir, el factor
$(\Lambda_\sst')^{t-1}$ del lema \ref{lema:fake}.
\end{itemize}

\paragraph{La acumulación} El paso 4 aplica $\J_{t,s}=\J_{t-1,s-1}+\F_{t,s}$ y da
\begin{equation}
\J^{A,r} =
\begin{pmatrix}
35.834 & 0.636 & 0.611 & 0.589 & 0.568 & 0.548 \\
35.550 & 36.457 & 1.236 & 1.190 & 1.147 & 1.107 \\
35.268 & 36.161 & 37.047 & 1.804 & 1.739 & 1.677 \\
34.987 & 35.868 & 36.741 & 37.606 & 2.344 & 2.261 \\
34.709 & 35.577 & 36.438 & 37.290 & 38.137 & 2.858 \\
34.432 & 35.289 & 36.138 & 36.979 & 37.813 & 38.643
\end{pmatrix}
\end{equation}

Ahora la matriz se ve completamente distinta, y la diferencia es instructiva.

\begin{itemize}[leftmargin=1.4em]
\item \textbf{La primera fila y la primera columna no cambiaron}: la recursión no
tiene de dónde acumular en los bordes. Esto se verifica en el código.
\item \textbf{Apareció una ``escalera'' en el triángulo inferior}, con valores
alrededor de $35$--$38$, que no estaba en $\F$. Son las columnas de $\F$
desplazadas y sumadas. Económicamente: una vez que el choque \emph{ocurrió}, el
efecto riqueza se acumuló y persiste.
\item \textbf{El triángulo superior sigue siendo pequeño} (de $0.5$ a $2.9$),
pero ahora crece al alejarse de la diagonal hacia abajo. Es la rama ascendente de
la ``carpa'' de la figura \ref{fig:columnas}: la acumulación previa al choque.
\end{itemize}

\paragraph{Verificar la recursión a mano} La proposición \ref{prop:principal}
dice que $\J_{2,2} = \F_{2,2}+\F_{1,1}+\F_{0,0}$. Con los números de arriba:
\begin{equation}
0.5895 + 0.6231 + 35.8340 = 37.0466 = \J_{2,2}. \quad\checkmark
\end{equation}
Y para una entrada fuera de la diagonal,
$\J_{4,2} = \F_{4,2}+\F_{3,1}+\F_{2,0} = 36.4377$, que es exactamente la entrada
correspondiente. Hacer esta verificación a mano con una matriz pequeña, la
primera vez que uno implementa el algoritmo, ahorra mucho tiempo después.

\begin{trampa}
Con $T=6$ el jacobiano es \emph{numéricamente} correcto pero
\emph{económicamente} inservible: supone que la economía vuelve al estado
estacionario en seis trimestres, lo que es falso. Las entradas de la esquina
inferior derecha ($38.6$) no se parecen a las del jacobiano verdadero con
$T=300$, porque allí el horizonte corto impone una condición terminal falsa. El
ejemplo sirve para ver el mecanismo, no para hacer economía.
\end{trampa}

\section{Diferenciación numérica}

En la práctica los pasos 1 y 2 se implementan con diferenciación numérica. Hay
tres opciones, en orden creciente de precisión y de esfuerzo de implementación:

\begin{enumerate}[leftmargin=1.4em]
\item \textbf{Un lado:} $(f(x+h)-f(x))/h$. Simple, error $O(h)$.
\item \textbf{Dos lados:} $(f(x+h)-f(x-h))/(2h)$. Cuesta el doble, error
$O(h^2)$, gana entre uno y dos órdenes de magnitud. Es lo que usa el código de
este manual.
\item \textbf{Diferenciación automática:} error cero, pero exige escribir el
paso hacia atrás con una librería de diferenciación automática.
\end{enumerate}

\begin{trampa}
Con diferenciación de un lado conviene usar \emph{corridas fantasma}: en lugar de
restar el valor de estado estacionario, correr el mismo código con $dx=0$ y
restar \emph{ese} resultado. Si el estado estacionario no convergió del todo, la
corrida fantasma arrastra el mismo sesgo que la perturbada y la resta lo cancela.
Es una línea de código y elimina una familia entera de errores.
\end{trampa}
